\chaptertitle{附录\quad 各章练习参考答案}

\sectiontitle{第一章\quad 平面上的几何变换}

\noindent\textbf{1.2 旋转}\\
1. $(0,3)$；$(0,-3)$ \quad 2. $(-2,-1)$ \quad 3. $(-3,4)$\\
4. $(\frac{\sqrt2}{2},\frac{\sqrt2}{2})$\\
5. $(1,1)\xrightarrow{90^\circ}(-1,1)\xrightarrow{90^\circ}(-1,-1)\xrightarrow{90^\circ}(1,-1)\xrightarrow{90^\circ}(1,1)$\\
6. $\theta=-90^\circ$：$\cos=0,\sin=-1$，$x'=y$，$y'=-x$。

\noindent\textbf{1.3 反射}\\
1. $(5,3)$；$(-5,-3)$ \quad 2. $(-2,4)$\\
3. 关于$y=x$：$(4,1)\to(1,4)$；关于$y=-x$：$(4,1)\to(-1,-4)$。\\
4. $(2,3)\to(2,-3)\to(-2,-3)$。等于关于原点对称（旋转$180^\circ$）。\\
5. $(1,0)\to(0,1)\to(0,-1)$。

\noindent\textbf{1.4 伸缩}\\
1. $(12,5)$；$(4,2.5)$ \quad 2. $(2,5)$，等价关于$y$轴对称。 \quad 3. $(6,4)$；$12$倍。\\
4. 一样，都得到$(2,3)$。 \quad 5. 能，$k=\frac12$即可。

\noindent\textbf{1.5 投影}\\
1. $(4,0)$；$(0,-3)$ \quad 2. 都重合到$(0,5)$；都重合到$(2,0)$。\\
3. 不同点可投影到同一点，无法唯一还原。

\noindent\textbf{1.6 切变}\\
1. $k=3$时：$(1,2)\to(7,2)$；$(1,0)\to(1,0)$不变因为$y=0$。\\
2. $(1,3)\to(-5,3)$，负$k$向左推。\\
3. $(0,0),(1,0),(4,1),(3,1)$。\\
4. $k=-2$可恢复：复合矩阵$\begin{pmatrix}1&2\\0&1\end{pmatrix}\begin{pmatrix}1&-2\\0&1\end{pmatrix}=I$。\\
5. $x'=x$，$y'=y+kx$。

\noindent\textbf{1.7 综合练习}\\
1. $(3,4)\xrightarrow{90^\circ}(-4,3)\xrightarrow{x\text{轴}}(-4,-3)$。\\
2. $(-1,5)\xrightarrow{y=x}(5,-1)\xrightarrow{x\text{拉伸}2}(10,-1)$。\\
3. 切变$(5,0)\to(5,0)$，$y$轴投影$(5,0)\to(0,0)$。\\
4. (a) $(0,0),(2,0),(2,-2),(0,-2)$\\
(b) $(0,0),(2,0),(2,-4),(0,-4)$\\
(c) $(-2,-4),(-4,-4)$（$k=1$切变）\\
(d) 旋转$180^\circ$取负：$(2,4),(4,4)$\\
5. (a)正确。旋转不改变$r$。 (b)错，不可逆。 (c)正确，底和高不变。 (d)正确。\\
6. 伸缩；切变。\\
7. 直线参数方程的两分量是$t$的一次函数，消$t$仍为直线。若含$x^2$或$\sqrt{x}$则会变为曲线。

\noindent\textbf{课后练习 A组}\\
1. $(2,7)$\\
2. $(-4,-6)$；$(6,-4)$\\
3. 拉伸后$(6,0)$，$y$轴反射后$(-6,0)$。\\
4. $k=-2$，$(3,-1)\to(3+(-2)(-1),-1)=(5,-1)$。向$x$轴投影$(5,0)$。\\
5. $90^\circ\times3$即$270^\circ$：$(0,0),(0,-1),(1,-1),(1,0)$。\\
6. $(2,5)\xrightarrow{180^\circ}(-2,-5)\xrightarrow{x\text{轴}}(-2,5)\xrightarrow{x\text{拉}3}(-6,5)$。\\
7. (a) $x'=-y,y'=x$\quad (b) $x'=-y,y'=-x$\quad (c) $x'=x/2,y'=3y$\\
8. 设$P=(x,y)$，旋转$90^\circ$：$x'=-y=2$，$y'=x=-3$。解得$x=-3,y=-2$，即$P=(-3,-2)$。\\
9. 旋转可逆，投影不可逆（信息丢失），切变可逆（做反向切变）。\\
10. 一次反射。三次=两次抵消+一次，等于一次反射。\\

\noindent\textbf{B组}\\
1. $x'=1\cdot\cos30^\circ-\sqrt3\cdot\sin30^\circ=\frac{\sqrt3}{2}-\frac{\sqrt3}{2}=0$；$y'=1\cdot\sin30^\circ+\sqrt3\cdot\cos30^\circ=\frac12+\frac32=2$。结果$(0,2)$。\\
2. $x'=4-x$，$y'=y$。\\
3. $\cos(\theta+(-\theta))=\cos0=1$，$\sin(\theta+(-\theta))=\sin0=0$，系数矩阵为$I$。\\
4. 总效果是旋转$\theta_1+\theta_2$。\\
5. 两个关于相交直线的反射=一个旋转。如关于$x$轴反射再关于$45^\circ$线反射=旋转$90^\circ$。\\
6. 由$(1,0)\to(0,2)$和$(0,1)\to(-3,0)$知$x'=-3y,y'=2x$。按成分：先旋转$90^\circ$得$(-y,x)$，再拉伸$x$方向$3$倍和$y$方向$2$倍。是旋转+伸缩的复合。\\
7. (a) 旋转$135^\circ$：$x'=-\frac{\sqrt2}{2}x-\frac{\sqrt2}{2}y,\;y'=\frac{\sqrt2}{2}x-\frac{\sqrt2}{2}y$。$(3,5)\to(-\frac{8\sqrt2}{2},-\frac{2\sqrt2}{2})=(-4\sqrt2,-\sqrt2)$。\\
(b) 关于$y=x$反射：$(-4\sqrt2,-\sqrt2)\to(-\sqrt2,-4\sqrt2)$。\\
(c) $x$方向拉伸$\frac13$：$(-\frac{\sqrt2}{3},-4\sqrt2)$。\\
8. 切变$k=1.5$矩阵$\begin{pmatrix}1&1.5\\0&1\end{pmatrix}$，四顶点变$(0,0),(2,0),(5,2),(3,2)$。再旋转$45^\circ$：用$\frac{\sqrt2}{2}\begin{pmatrix}1&-1\\1&1\end{pmatrix}$作用得$(0,0),(\sqrt2,\sqrt2),(\frac{3\sqrt2}{2},\frac{7\sqrt2}{2}),(\frac{\sqrt2}{2},\frac{5\sqrt2}{2})$。\\
9. 反推：先撤销反射（关于$y=-x$反射$=$自身），再撤销$y$方向拉伸$3$倍（压缩到$\frac13$倍）。关于$y=-x$反射$(5,-2)\to(2,-5)$。再$y$压缩$\frac13$倍：$(2,-\frac53)$。原坐标为$(2,-\frac53)$。\\
10. 例如先$x$方向拉伸$2$倍再旋转$45^\circ$（既非纯伸缩也非纯旋转）。公式：$\frac{\sqrt2}{2}\begin{pmatrix}1&-1\\1&1\end{pmatrix}\begin{pmatrix}2&0\\0&1\end{pmatrix}=\frac{\sqrt2}{2}\begin{pmatrix}2&-1\\2&1\end{pmatrix}$，即$x'=\sqrt2\,x-\frac{\sqrt2}{2}y,\;y'=\sqrt2\,x+\frac{\sqrt2}{2}y$。验证$(1,0)\to(\sqrt2,\sqrt2)$。

\sectiontitle{第二章\quad 矩阵——变换的代码}

\noindent\textbf{2.1 矩阵乘向量}\\
1. $\begin{pmatrix}0\\5\end{pmatrix}$ \quad 2. $\begin{pmatrix}-3\\-2\end{pmatrix}$ \quad 3. $\begin{pmatrix}4\\5\end{pmatrix}$ \quad 4. $\begin{pmatrix}5\\1\end{pmatrix}$\\

\noindent\textbf{2.2 矩阵乘法}\\
1. $\begin{pmatrix}5&8\\9&12\end{pmatrix}$ \quad 2. $\begin{pmatrix}2&7\\6&15\end{pmatrix}$，不同。\\
3. $\begin{pmatrix}-1&0\\0&1\end{pmatrix}$\quad 4. $\begin{pmatrix}-1&0\\0&1\end{pmatrix}$\\
5. $\begin{pmatrix}kp&0\\0&lq\end{pmatrix}$，结果相同。\quad 6. $\begin{pmatrix}0&0\\0&0\end{pmatrix}$，零矩阵。\\

\noindent\textbf{2.3 特殊矩阵}\\
1. $1$ \quad 2. $k,l\neq0$时可逆。\quad 3. 旋转矩阵如$\begin{pmatrix}\cos\theta&-\sin\theta\\\sin\theta&\cos\theta\end{pmatrix}$。\\

\noindent\textbf{课后练习 A组}\\
1. $\begin{pmatrix}3\cdot1+1\cdot4&3\cdot0+1\cdot2\\2\cdot1+5\cdot4&2\cdot0+5\cdot2\end{pmatrix}=\begin{pmatrix}7&2\\22&10\end{pmatrix}$\\
2. $\begin{pmatrix}3&1\\14&31\end{pmatrix}$，不同。\\
3. $\det=2\times3-5\times1=1$。$A^{-1}=\begin{pmatrix}3&-5\\-1&2\end{pmatrix}$。\\
4. 增广矩阵$\begin{pmatrix}3&-2&|&8\\1&4&|&2\end{pmatrix}\xrightarrow{R_1\leftrightarrow R_2}\begin{pmatrix}1&4&|&2\\3&-2&|&8\end{pmatrix}\xrightarrow{R_2-3R_1}\begin{pmatrix}1&4&|&2\\0&-14&|&2\end{pmatrix}$。回代$y=-\frac17$，$x=2-4y=\frac{18}{7}$。\\
5. (a) $\begin{pmatrix}-\frac12&-\frac{\sqrt3}{2}\\\frac{\sqrt3}{2}&-\frac12\end{pmatrix}$\\
(b) $R_{90}S_{y=-x}=\begin{pmatrix}0&-1\\1&0\end{pmatrix}\begin{pmatrix}0&-1\\-1&0\end{pmatrix}=\begin{pmatrix}1&0\\0&-1\end{pmatrix}$（关于$x$轴反射）。\\
6. $R_{90}S_{x2}=\begin{pmatrix}0&-1\\1&0\end{pmatrix}\begin{pmatrix}2&0\\0&1\end{pmatrix}=\begin{pmatrix}0&-1\\2&0\end{pmatrix}$。交换：$\begin{pmatrix}2&0\\0&1\end{pmatrix}\begin{pmatrix}0&-1\\1&0\end{pmatrix}=\begin{pmatrix}0&-2\\1&0\end{pmatrix}$。\\
7. $A^2=\begin{pmatrix}0&1\\1&0\end{pmatrix}^2=\begin{pmatrix}1&0\\0&1\end{pmatrix}=I$。两次关于$y=x$反射回到原位。\\
8. 代入任意数值，$(AB)C=A(BC)$恒成立。矩阵乘法满足结合律。\\
9. 增广矩阵$\begin{pmatrix}2&1&-1&|&3\\1&-2&3&|&1\\3&1&2&|&7\end{pmatrix}\xrightarrow{R_1\leftrightarrow R_2}\begin{pmatrix}1&-2&3&|&1\\2&1&-1&|&3\\3&1&2&|&7\end{pmatrix}\xrightarrow{R_2-2R_1,R_3-3R_1}\begin{pmatrix}1&-2&3&|&1\\0&5&-7&|&1\\0&7&-7&|&4\end{pmatrix}\xrightarrow{R_3-R_2}\begin{pmatrix}1&-2&3&|&1\\0&5&-7&|&1\\0&2&0&|&3\end{pmatrix}$回代$x=1,y=3/2,z=1/2$。\\
10. $\det(AB)=2\times0.5=1$。$\det(A^{-1})=1/2$。\\

\noindent\textbf{B组}\\
1. 由$AB=I$知$B$是$A$的右逆。对方阵若右逆存在则左逆也存在且相等，故$BA=I$。证明需用到行列式非零。\\
2. 所有反射矩阵（关于任何过原点直线的反射）和$\pm I$以及$\begin{pmatrix}\pm1&0\\0&\pm1\end{pmatrix}$中使得平方为$I$的组合。\\
3. 增广矩阵化为$\begin{pmatrix}1&a&|&1\\0&1-a^2&|&1-a\end{pmatrix}$。$a\neq\pm1$时唯一解，$a=1$时无穷多解，$a=-1$时无解。\\
4. 复合结果为另一个反射矩阵。旋转加反射=反射，因为行列式为$-1$且正交。\\
5. $ad-bc=0$意味着两列向量共线，变换把平面压扁成一条线，无数点映射到同一点，无法找到逆映射。\\
6. 增广矩阵$\begin{pmatrix}3&2&|&5\\\frac12&-\frac13&|&1\end{pmatrix}\xrightarrow{R_2- \frac16 R_1}\begin{pmatrix}3&2&|&5\\0&-\frac23&|&\frac16\end{pmatrix}$。回代$y=-\frac14$，$3x+2(-\frac14)=5$，$x=\frac{11}{6}$。\\
7. $A^{-1}=\begin{pmatrix}\frac12&-\frac{\sqrt3}{2}\\\frac{\sqrt3}{2}&\frac12\end{pmatrix}$——这正是旋转$-60^\circ$的矩阵。逆矩阵对应逆向旋转，与$A$的几何含义（旋转$60^\circ$）互为逆。\\
8. $AB=\begin{pmatrix}\frac{11}{15}&\frac{2}{15}\\-\frac{2}{15}&\frac{11}{15}\end{pmatrix}$。$A$不是正交矩阵（$A^{\text T}A=\frac59 I$，它是一个旋转矩阵乘以$\frac{\sqrt5}{3}$），故$AB$也不是正交矩阵。若两个因子均为正交矩阵，乘积才必为正交矩阵。\\
9. 不可逆（$\det=0$）。$(x,y,z)\to(2x,3y,0)$。所有$z$坐标不同的点只要$(x,y)$相同就会被压缩到同一点——例如$(1,1,0)$和$(1,1,5)$都映射为$(2,3,0)$。

\sectiontitle{第三章\quad 行列式}

\noindent\textbf{3.2 负行列式}\\
1. $\det=-1$，负。 \quad 2. $\det=1$，正。$180^\circ$旋转不反转定向——虽然每个向量都反向，但手性（从$x$轴到$y$轴的旋转方向）不变。\\
3. 如$\begin{pmatrix}0&1\\1&0\end{pmatrix}$（关于$y=x$反射）。

\noindent\textbf{3.3 行列式为0}\\
1. $\det=3\times4-6\times2=0$，共线，线性相关。 \quad 2. $\det=0$，列向量$(1,0)$和$(2,0)$线性相关，平面被压成$x$轴。\\
3. 不共线。例如$\begin{pmatrix}2&3\\1&4\end{pmatrix}$。\quad 4. $\det=1\neq0$，可逆。

\noindent\textbf{3.4 三阶行列式}\\
1. 行列式为$0$，不可逆。\quad 2. 行列式为$0$（两列相等意味着线性相关）。

\noindent\textbf{课后练习 A组}\\
1. (a) $4\times2-3\times1=5$，放大$5$倍。(b) $35$，放大$35$倍。(c) $-6$，面积放大$6$倍且定向反转。\\
2. (a) $\det=0$，$(1,2)$和$(2,4)=2(1,2)$成比例。(b) $\det=-6\neq0$。(c) $\det=1\neq0$。\\
3. $1\cdot3\cdot4+2\cdot1\cdot2+0\cdot0\cdot0-0\cdot3\cdot2-2\cdot0\cdot4-1\cdot1\cdot0=12+4-0-0-0=16$。\\
4. $\det(2A)=2^2\det A=4\times3=12$（每行翻倍，两行各翻倍）。\\
5. $\det(AB)=\det A\cdot\det B=-10$，$\det(BA)=-10$（相同）。\\
6. $\det=4-6=-2\neq0$，可逆。$\det(A^{-1})=-1/2$。\\
7. $\det=x^2-1$。$x=\pm1$时行列式为$0$，此时两列向量成比例：当$x=1$时$\begin{pmatrix}1\\1\end{pmatrix}$和$\begin{pmatrix}1\\1\end{pmatrix}$相等；当$x=-1$时互为相反。几何上平面被压成一条线。\\
8. $0$（一旦有一行全为零，行列式展开式中每一项都含该行的零因子）。\\
9. $\det=d_1d_2d_3$。任一$d_i=0$时不可逆。\\
10. 例如$\begin{pmatrix}1&2\\3&6\end{pmatrix}$（沿$(1,3)$方向）和$\begin{pmatrix}2&3\\4&6\end{pmatrix}$（沿$(2,4)$方向）。\\

\noindent\textbf{B组}\\
1. 不变。属于"某列加另一列的倍数"性质。\\
2. $\det(A^2)=\det(A\cdot A)=\det A\cdot\det A=(\det A)^2$。反例：$\det A\neq1$时$\det(A^2)\neq\det A$，如$\det A=2$则$\det(A^2)=4\neq2$。\\
3. 展开后为$a^3+b^3+c^3-3abc=(a+b+c)(a^2+b^2+c^2-ab-bc-ca)$。\\
4. 由$A^{\text T}=-A$，取行列式：$\det(A^{\text T})=\det(-A)=(-1)^n\det A$。又$\det(A^{\text T})=\det A$，得$\det A=(-1)^n\det A$。若$n$为奇数，$\det A=-\det A$，故$\det A=0$。\\
5. $\det\begin{pmatrix}3&7\\x&y\end{pmatrix}=\det\begin{pmatrix}2+1&3+4\\x&y\end{pmatrix}=1+2=3$。\\
6. 通过行变换：$R_2\leftarrow R_2-\frac12 R_1$，$R_3\leftarrow R_3+R_1$，再$R_3\leftarrow R_3-\frac12 R_2$，化为上三角$\begin{pmatrix}2&-1&3\\0&\frac12&\frac52\\0&0&3\end{pmatrix}$，$\det=2\cdot\frac12\cdot3=3$。\\
7. $\det A=k^2-6=0$，$k=\pm\sqrt6$。此时两列成比例：$\begin{pmatrix}k\\3\end{pmatrix}$与$\begin{pmatrix}2\\k\end{pmatrix}$满足$k/2=3/k$。\\
8. 通分后计算：$\det=\frac{1}{2160}$（先提取公分母$60$的倒数，再展开三阶行列式）。\\
9. 行列式为$0$。因为第三行是第一行和第二行的线性组合，三个行向量线性相关。\\

\sectiontitle{第四章\quad 特征值}

\noindent\textbf{4.1 观察特征方向}\\
1. $(1,0)$是特征向量，特征值$4$；$(0,1)$也是特征向量，特征值$0.5$。\\
2. 沿$x$轴方向的向量$(t,0)$切变后仍是$(t,0)$，是特征向量（特征值$1$）。沿$y$轴方向的向量切变后方向改变，不是特征向量。\\
3. 几何上：旋转改变了所有非零向量的方向。代数上：特征方程$\lambda^2+1=0$无实数根。

\noindent\textbf{4.2 特征值与特征向量}\\
1. $\lambda_1=5,v_1=(1,0)$；$\lambda_2=2,v_2=(0,1)$。\\
2. $\lambda_1=3,v_1=(1,1)$；$\lambda_2=-1,v_2=(1,-1)$。\\
3. $\lambda=-1$（二重），任何非零向量。\\
4. 特征方程$\lambda^2-5\lambda+6=0$，$\lambda=2,3$。

\noindent\textbf{4.3 旋转与实特征向量}\\
1. 旋转$60^\circ$矩阵$A=\begin{pmatrix}\frac12&-\frac{\sqrt3}{2}\\\frac{\sqrt3}{2}&\frac12\end{pmatrix}$。特征方程$\lambda^2-\lambda+1=0$，判别式$1-4=-3<0$，无实数根。\\
2. $\lambda^2+4=0$的解为$\lambda=\pm2i$，说明矩阵含有旋转$90^\circ$的成分。

\noindent\textbf{4.4 特征值的和与积}\\
1. $\lambda^2-1\lambda-6=0$，$\lambda_1=3,\lambda_2=-2$。\\
2. $\det A=0$时至少有一个特征值为$0$，另一个不能唯一确定。若$\tr A=4$，则$\lambda_1=0,\lambda_2=4$。

\noindent\textbf{4.5 施密特正交化}\\
1. $\vec u_1=(3,0)$，$\vec u_2=(0,4)$。原向量已互相垂直，正交化不改变它们。\\
2. $\vec u_1=(1,2)$。投影系数$=\frac{3+8}{1+4}=\frac{11}{5}$。$\vec u_2=(3,4)-\frac{11}{5}(1,2)=(\frac{4}{5},-\frac{2}{5})$。验证：$(1,2)\cdot(\frac{4}{5},-\frac{2}{5})=0$。

\noindent\textbf{4.6 对角化}\\
1. $\vec u_1=(2,1),\vec u_2=(-1,2)$。\quad 2. $\lambda=3,v=(1,0)$；$\lambda=5,v=(0,1)$。\quad 3. 一个或两个，取决于迹。

\noindent\textbf{课后练习 A组}\\
1. (a) $\lambda^2-7\lambda+10=0$，$\lambda_1=5,\lambda_2=2$。$\lambda=5$：$v=(1,-1)$？——求解$(A-5I)v=\begin{pmatrix}-1&1\\2&-2\end{pmatrix}v=0$得$v_1=(1,1)$。$\lambda=2$：$v_2=(1,-2)$。\\
(b) $\lambda=3$（二重），任意向量。\\
(c) $\lambda=1,v=(1,1)$；$\lambda=-1,v=(1,-1)$。\\
2. $\lambda^2-5\lambda+4=0$，$\lambda=1,4$。$v_{\lambda=1}=(1,-1)$，$v_{\lambda=4}=(1,2)$。\\
3. 特征方程$\lambda^2-6\lambda+8=0$，$\lambda=2,4$。\\
4. $\lambda=3$（二重）。特征向量只有$(1,0)$一个方向。几何上是一个切变加等比例伸缩，不可对角化。\\
5. 用$A=PDP^{-1}$构造：$P=\begin{pmatrix}1&1\\0&1\end{pmatrix}$，$D=\begin{pmatrix}2&0\\0&-1\end{pmatrix}$。$A=PDP^{-1}=\begin{pmatrix}1&1\\0&1\end{pmatrix}\begin{pmatrix}2&0\\0&-1\end{pmatrix}\begin{pmatrix}1&-1\\0&1\end{pmatrix}=\begin{pmatrix}2&-3\\0&-1\end{pmatrix}$。验证：特征方程$(\lambda-2)(\lambda+1)=0$，$\lambda=2,-1$。\\
6. $\lambda=2$（二重），特征向量$(1,0)$唯一。每个点被拉长两倍并附加切变。\\
7. $\vec u_1=(3,0),\vec u_2=(0,2)$（巧合，原向量已正交）。\\
8. 需$\lambda_1+\lambda_2=5,\lambda_1\lambda_2=6$。解为$\lambda=2,3$。可取$A=\begin{pmatrix}2&0\\0&3\end{pmatrix}$。\\

\noindent\textbf{B组}\\
1. $A^{-1}$的特征值为$1/\lambda_i$，特征向量相同。验证：$A v=\lambda v\Rightarrow v=\lambda A^{-1}v\Rightarrow A^{-1}v=\frac1\lambda v$。\\
2. $A^2 v=A(Av)=A(\lambda v)=\lambda Av=\lambda^2 v$。例：对角阵$\begin{pmatrix}2&0\\0&3\end{pmatrix}^2=\begin{pmatrix}4&0\\0&9\end{pmatrix}$，特征值$4,9=2^2,3^2$。\\
3. 由$A v=\lambda v$，左乘$A^{-1}$得$v=\lambda A^{-1}v\Rightarrow A^{-1}v=\frac1\lambda v$，特征向量相同。\\
4. $(A^2-2A+I)v=(A^2)v-2Av+Iv=(\lambda^2-2\lambda+1)v$。$\lambda=2\Rightarrow$特征值$1$；$\lambda=-3\Rightarrow$特征值$16$。\\
5. $A=\begin{pmatrix}0&a\\-a&0\end{pmatrix}$，特征方程$\lambda^2+a^2=0$，$\lambda=\pm ai$，无实特征值。\\
6. 特征方程$\lambda^2-4\lambda-5=0$，$\lambda_1=5,\lambda_2=-1$。$v_1=(1,1)$，$v_2=(-2,1)$。$v_1\cdot v_2=-1\neq0$，不正交。施密特：$\vec u_1=(1,1)$，$\vec u_2=(-2,1)-\frac{-1}{2}(1,1)=(-\frac32,\frac32)$。\\
7. 特征方程$\lambda^2-3\lambda+2=0$，$\lambda=1,2$，均为实数（因为$\tr A=3,\det A=2$，韦达定理保证有实根）。矩阵不含纯旋转成分，但含切变+伸缩。\\
8. $\lambda_1+\lambda_2=0,\lambda_1\lambda_2=-2$，得$\lambda_1=\sqrt2,\lambda_2=-\sqrt2$。不可能是旋转——旋转矩阵的行列式为$+1$而非$-2$，且特征值非纯虚数，说明矩阵含伸缩成分而非纯旋转。\\
9. $\vec u_1=(1,3)$。投影系数$=\frac{2-3}{1+9}=-\frac1{10}$。$\vec u_2=(2,-1)+\frac1{10}(1,3)=(\frac{21}{10},-\frac{7}{10})$。验证$(1,3)\cdot(\frac{21}{10},-\frac{7}{10})=\frac{21}{10}-\frac{21}{10}=0$。方向变化是因为原始$\vec v_1,\vec v_2$不正交，正交化过程中$\vec v_2$被减去了沿$\vec v_1$的分量。A组第7题中$(3,0)$和$(1,2)$的施密特结果$(3,0),(0,2)$恰与原向量同向——这是因为$(3,0)$与$(1,2)$的点积不为零，但减去的投影恰好将$x$分量清零，属于数值巧合。

\sectiontitle{第五章\quad 齐次坐标与仿射变换}

\noindent\textbf{5.1 平移的局限}\\
1. (a) 旋转$90^\circ$保持原点； (b) 平移$(2,-1)$不保持原点（原点变为$(2,-1)$）； (c) 关于$x$轴反射保持原点； (d) 先平移再旋转，原点变为$(0+3,0+0)\xrightarrow{180^\circ}(-3,0)$，不保持原点。\\
2. 因为$x'=ax+by$和$y'=cx+dy$在$(x,y)=(0,0)$时必然得$(0,0)$。即线性变换恒将原点固定，无法移动原点。

\noindent\textbf{5.2 齐次坐标}\\
1. $R_{90}$的齐次矩阵$\begin{pmatrix}0&-1&0\\1&0&0\\0&0&1\end{pmatrix}$，平移$(2,0)$的矩阵$\begin{pmatrix}1&0&2\\0&1&0\\0&0&1\end{pmatrix}$。复合$=T R_{90}=\begin{pmatrix}0&-1&2\\1&0&0\\0&0&1\end{pmatrix}$。$(1,3)\to(-1,3)$。\\
2. $\begin{pmatrix}1&0&a\\0&1&b\\0&0&1\end{pmatrix}$。\\
3. 若第三行不是$(0,0,1)$，则第三个坐标不再是恒为$1$的常量，不再对应平面上的仿射变换（可能表示投影变换或透视变换）。

\noindent\textbf{5.4 变换的分解与复合}\\
1. 先旋转$90^\circ$再平移$(2,0)$：$T_{\text{平移}}R_{90}=\begin{pmatrix}1&0&2\\0&1&0\\0&0&1\end{pmatrix}\begin{pmatrix}0&-1&0\\1&0&0\\0&0&1\end{pmatrix}=\begin{pmatrix}0&-1&2\\1&0&0\\0&0&1\end{pmatrix}$。原点$(0,0)\to(2,0)$。

先平移$(2,0)$再旋转$90^\circ$：$R_{90}T_{\text{平移}}=\begin{pmatrix}0&-1&0\\1&0&0\\0&0&1\end{pmatrix}\begin{pmatrix}1&0&2\\0&1&0\\0&0&1\end{pmatrix}=\begin{pmatrix}0&-1&0\\1&0&2\\0&0&1\end{pmatrix}$。原点$(0,0)\to(0,2)$。

结果不同——因为旋转改变了平移的方向。先平移再旋转，平移量被旋转了。\\

\noindent\textbf{课后练习 A组}\\
1. $\begin{pmatrix}1&0&1\\0&1&-2\\0&0&1\end{pmatrix}\begin{pmatrix}2\\3\\1\end{pmatrix}=\begin{pmatrix}3\\1\\1\end{pmatrix}$，即$(3,1)$。\\
2. (a) $\begin{pmatrix}\frac{\sqrt2}{2}&-\frac{\sqrt2}{2}&5\\\frac{\sqrt2}{2}&\frac{\sqrt2}{2}&0\\0&0&1\end{pmatrix}$\quad (b) $\begin{pmatrix}1&0&0\\0&0&-3\\0&0&1\end{pmatrix}$\\
3. $\begin{pmatrix}5\\3\\1\end{pmatrix}$即$(5,3)$。线性部分$\det=3\times2=6$，面积$6$倍。\\
4. 旋转+平移保持面积（$|\det|=1$）；投影+平移降维不保持面积（$\det=0$）；切变+平移保持面积（$\det=1$）；反射+平移保持绝对面积（$|\det|=1$）。\\
5. 矩阵$\begin{pmatrix}1&0&0\\0&-1&-2\\0&0&1\end{pmatrix}$，$(3,4)\to(3,-6)$。\\

\noindent\textbf{B组}\\
1. 由$(0,0)\to(1,2)$得平移$(1,2)$。$(1,0)\to(4,2)$得$x'=3x+1$即$a=3$。$(0,1)\to(1,5)$得$y'=3y+2$即$d=3$。矩阵为$\begin{pmatrix}3&0&1\\0&3&2\\0&0&1\end{pmatrix}$。\\
2. $4\times4$矩阵，形式为$\begin{pmatrix}1&0&0&a\\0&1&0&b\\0&0&1&c\\0&0&0&1\end{pmatrix}$。\\
3. 两个旋转+平移 = 新的旋转+新的平移。新旋转=两个旋转矩阵的乘积（仍是正交矩阵），新平移=第一个平移旋转后加上第二个平移。\\
4. 由$(0,0)\to(2,-1)$得$\vec t=(2,-1)$。$(1,0)\to(5,1)$得$a+2=5,c-1=1\Rightarrow a=3,c=2$。$(0,1)\to(0,3)$得$b+2=0,d-1=3\Rightarrow b=-2,d=4$。矩阵$\begin{pmatrix}3&-2&2\\2&4&-1\\0&0&1\end{pmatrix}$。分解：线性部分$A=\begin{pmatrix}3&-2\\2&4\end{pmatrix}$，平移$\vec t=(2,-1)$。\\
5. 先旋转后平移：$T R_{60}=\begin{pmatrix}1&0&4\\0&1&0\\0&0&1\end{pmatrix}\begin{pmatrix}\frac12&-\frac{\sqrt3}{2}&0\\\frac{\sqrt3}{2}&\frac12&0\\0&0&1\end{pmatrix}$，$(1,1)\to(\frac12-\frac{\sqrt3}{2}+4,\frac{\sqrt3}{2}+\frac12)$。先平移后旋转：$R_{60}T$，$(1,1)\to(\frac12(1+4)-\frac{\sqrt3}{2}\cdot1,\frac{\sqrt3}{2}(1+4)+\frac12\cdot1)=(\frac52-\frac{\sqrt3}{2},\frac{5\sqrt3}{2}+\frac12)$。结果不同。\\
6. 中心为$(1.5,1)$，绕中心旋转$180^\circ$等价于绕原点旋转$180^\circ$再加平移修正。齐次矩阵$\begin{pmatrix}-1&0&8\\0&-1&2\\0&0&1\end{pmatrix}$。四角点$(0,0)\to(8,2)$，$(3,0)\to(5,2)$，$(3,2)\to(5,0)$，$(0,2)\to(8,0)$。

\sectiontitle{第六章\quad 二次型}

\noindent\textbf{6.1 交叉项识别}\\
1. (a) 无$xy$项；(b) 有$-2xy$项；(c) 无$xy$项；(d) 有$3xy$项。\\
2. (a) $2x^2+3y^2=5$，即$\frac{x^2}{2.5}+\frac{y^2}{5/3}=1$，椭圆。(c) $4x^2-y^2=1$，双曲线。

\noindent\textbf{6.2 二次型的矩阵}\\
1. (a) $M=\begin{pmatrix}1&3\\3&9\end{pmatrix}$。(b) $M=\begin{pmatrix}2&-1\\-1&1\end{pmatrix}$。(c) $M=\begin{pmatrix}1&0\\0&1\end{pmatrix}$。\\
2. 代入验证：$(x,y)M\begin{pmatrix}x\\y\end{pmatrix}$展开即得原二次型。

\noindent\textbf{6.3 特征值消交叉项}\\
1. $M=\begin{pmatrix}3&1\\1&3\end{pmatrix}$，$\lambda^2-6\lambda+8=0$，$\lambda=4,2$，均正，椭圆。\\
2. 坐标旋转后，二次型化为$\lambda_1X^2+\lambda_2Y^2$，其中$\lambda_1,\lambda_2$恰好是$M$的特征值。因此直接求特征值等价于先旋转再分类——不需要实际计算旋转角。

\noindent\textbf{6.4 用特征值分类}\\
1. (a) $M=\begin{pmatrix}2&2\\2&5\end{pmatrix}$，$\det M=10-4=6>0$，椭圆。(b) $M=\begin{pmatrix}1&-3\\-3&9\end{pmatrix}$，$\det M=9-9=0$，抛物线。\\
2. $\det M=\lambda_1\lambda_2=0$说明至少有一个特征值为零。在标准形$\lambda_1X^2+\lambda_2Y^2$中有一项消失，方程退化为$Y^2=\text{常数}$或$X^2=\text{常数}$的形式——即抛物线（或退化为两条平行线/重合直线/一个点）。几何上，行列式为零意味着变换降维，圆锥曲线退化为"开口"形状。

\noindent\textbf{课后练习 A组}\\
1. (a) $\begin{pmatrix}1&1\\1&3\end{pmatrix}$\quad (b) $\begin{pmatrix}4&-3\\-3&1\end{pmatrix}$\\
2. (a) $M=\begin{pmatrix}3&1\\1&3\end{pmatrix}$，$\lambda=4,2$，均正，椭圆。\\
(b) $M=\begin{pmatrix}2&2\\2&5\end{pmatrix}$，$\lambda=6,1$，椭圆。\\
(c) $M=\begin{pmatrix}4&6\\6&9\end{pmatrix}$，特征方程$\lambda^2-13\lambda=0$，$\lambda=13,0$，抛物线。\\
3. $M=\begin{pmatrix}2&0\\0&3\end{pmatrix}$，特征值$2,3$。\\
4. 一正一负，双曲线。\\
5. $(x+2y)^2=0$，即直线$x+2y=0$。从特征值看：$M=\begin{pmatrix}1&2\\2&4\end{pmatrix}$，$\lambda=5,0$，退化抛物线（直线）。\\
6. 特征值一正一负，双曲线。\\

\noindent\textbf{B组}\\
1. $\det M=\lambda_1\lambda_2=ac-b^2$。这正是二次型的判别式乘以$-1$。判别式决定曲线类型与特征值正负号分类等价。\\
2. 当$a=c$时$\tan2\theta=\frac{2b}{0}$无定义，说明$\theta=45^\circ$。代入旋转公式可直接验证消去$xy$项。\\
3. 特征值$\lambda$对应平方项$\lambda X^2$。椭圆方程写为$\lambda_1X^2+\lambda_2Y^2=1$，即$\frac{X^2}{1/\lambda_1}+\frac{Y^2}{1/\lambda_2}=1$。所以半轴长为$1/\sqrt{\lambda_i}$。较小的特征值对应较长的半轴。\\
4. $x^2+2xy+y^2=(x+y)^2$。特征值$M=\begin{pmatrix}1&1\\1&1\end{pmatrix}$，$\lambda=2,0$，退化抛物线即一条直线。\\
5. $M=\begin{pmatrix}3&4\\4&-3\end{pmatrix}$，$\lambda^2-0\lambda-25=0$，$\lambda=5,-5$。一正一负，双曲线。$\lambda=5$：$v=(2,1)$；$\lambda=-5$：$v=(1,-2)$。实轴沿$(2,1)$，虚轴沿$(1,-2)$。\\
6. $M=\begin{pmatrix}6&-\frac{\sqrt3}{2}\\-\frac{\sqrt3}{2}&5\end{pmatrix}$，$\lambda^2-11\lambda+\frac{117}{4}=0$，$\lambda=\frac{11\pm2}{2}$即$\lambda_1=\frac{13}{2},\lambda_2=\frac{9}{2}$。均正，椭圆。\\
7. 特征值相等$\Rightarrow\lambda_1=\lambda_2$。由$\lambda_1+\lambda_2=a+c=2\lambda$和$\lambda_1\lambda_2=ac-b^2=\lambda^2$，得$(a-c)^2+4b^2=0\Rightarrow a=c,b=0$，即$M=a I$。此时二次型为$a(x^2+y^2)$，对应的$\lambda_1=\lambda_2=a$，曲线为圆（若$a>0$）或不存在（若$a<0$）。特征值相等但非圆的情况不可能出现，因为$M$对称且特征值相等$\Rightarrow M=\lambda I$。\\
8. $\alpha+\gamma=\tr M=7+(-2)=5$（给定），$\alpha\gamma-\beta^2=\det M=-14$。还需额外条件才能唯一确定$\alpha,\beta,\gamma$。若取$\alpha=7,\gamma=-2$，则$7(-2)-\beta^2=-14\Rightarrow\beta^2=0\Rightarrow\beta=0$。一组解为$\alpha=7,\beta=0,\gamma=-2$（曲线$7x^2-2y^2=1$，双曲线）。

\sectiontitle{第七章\quad 推广与展望}

\noindent\textbf{7.1 矩阵的推广}\\
1. 秩为$1$（三列中第二列为第一列的$2$倍，第三列为第一列的$3$倍，仅第一列线性无关）。\\
2. 秩最大为$3$（取$\min(3,5)=3$），最小为$0$（全零矩阵）。

\noindent\textbf{7.2 $n$阶行列式与拉普拉斯展开}\\
1. 按第一列展开，第一列非零元为$1$和$3$：\\
$\det = 1\cdot(-1)^{1+1}\begin{vmatrix}0&1\\2&4\end{vmatrix}+3\cdot(-1)^{2+1}\begin{vmatrix}2&0\\2&4\end{vmatrix}+0 = 1\cdot(0-2)-3\cdot(8-0) = -2-24 = -26$。\\
验证：按第三章公式$=1\cdot0\cdot4+2\cdot1\cdot0+0\cdot3\cdot2-0\cdot0\cdot0-2\cdot3\cdot4-1\cdot1\cdot2=0+0+0-0-24-2=-26$，一致。\\
2. 零元多意味着许多项本身就是$0$，不需要计算对应的余子式，大幅减少计算量。

\noindent\textbf{7.3 克拉默法则}\\
1. $\det A = 3\times4-(-2)\times1 = 14$。\\
$\det A_1 = \begin{vmatrix}8&-2\\2&4\end{vmatrix}=8\cdot4-(-2)\cdot2=36$，$x=36/14=18/7$。\\
$\det A_2 = \begin{vmatrix}3&8\\1&2\end{vmatrix}=3\cdot2-8\cdot1=-2$，$y=-2/14=-1/7$。\\
与高斯消元一致（见第二章答案A组第4题）。\\
2. 克拉默法则：每个$n$阶行列式约需$n!$次乘法，共需$(n+1)n!\approx n\cdot n!$次。$n=100$时约$100\times100!\approx9.3\times10^{159}$次——完全不可行。高斯消元约需$n^3=10^6$次。因此克拉默法则仅用于理论分析，不作实际数值计算。

\noindent\textbf{7.4 奇异值分解简介}\\
1. $A^{\text T}A = \begin{pmatrix}1&0&0\\0&2&0\end{pmatrix}\begin{pmatrix}1&0\\0&2\\0&0\end{pmatrix}=\begin{pmatrix}1&0\\0&4\end{pmatrix}$。特征值为$1$和$4$——$A$在水平方向放大$1$倍，在竖直方向放大$2$倍（对应的奇异值为$\sqrt{1}=1$和$\sqrt{4}=2$）。\\
2. 非零奇异值的个数$= \operatorname{rank}(A)=r$。因为$A^{\text T}A$的秩等于$A$的秩，而$A^{\text T}A$的非零特征值个数等于其秩，非零奇异值为这些特征值的平方根。

\noindent\textbf{7.5 全书回顾与综合练习}\\
1. (a) $\det A = 2\times2-1\times1=3$，面积放大$3$倍。\\
(b) 特征方程$\lambda^2-4\lambda+3=0$，$\lambda_1=3,\lambda_2=1$。$\lambda=1$：$v_1=(1,-1)$（沿此方向不变）；$\lambda=3$：$v_2=(1,1)$（沿此方向拉伸$3$倍）。\\
(c) $2x^2+2xy+2y^2$。$M=\begin{pmatrix}2&1\\1&2\end{pmatrix}$，$\lambda=3,1$均正，椭圆。\\
(d) $P=\begin{pmatrix}1&1\\-1&1\end{pmatrix}$，$D=\begin{pmatrix}1&0\\0&3\end{pmatrix}$。$P^{-1}=\frac12\begin{pmatrix}1&-1\\1&1\end{pmatrix}$。

2. (a) $A=\begin{pmatrix}2&1\\3&4\end{pmatrix}$。\\
(b) $\det A = 8-3=5$，面积放大$5$倍。\\
(c) $\det A=5\neq0$，可逆。$A^{-1}=\frac15\begin{pmatrix}4&-1\\-3&2\end{pmatrix}$。\\
(d) 特征方程$\lambda^2-6\lambda+5=0$，$\lambda_1=5,\lambda_2=1$。特征向量均存在，方向不变。

3. (a) $M=\begin{pmatrix}2&1\\1&2\end{pmatrix}$。\\
(b) $\lambda^2-4\lambda+3=0$，$\lambda=3,1$，均正，椭圆。\\
(c) $\lambda=3$：$v=(1,1)$（短轴）；$\lambda=1$：$v=(1,-1)$（长轴）。\\
(d) 标准方程：$3X^2+Y^2=3$，即$\frac{X^2}{1}+\frac{Y^2}{3}=1$。

4. (a) 三个列向量线性无关（因为$\det\neq0$）。\\
(b) $\det A = 1\times1\times1=1$。切变不改变三维体积。\\
(c) 特征方程$(1-\lambda)^3=0$，$\lambda=1$（三重）。$A$只有一个独立的特征向量$(1,0,0)$，不可对角化。

5.（开放题，要点）几何角度：$\det A=0$意味着变换把空间压扁到更低维度，多个不同的原像可能映射到同一个像，无法唯一还原——因此逆变换不存在。代数角度：$\det A=0$时$A\vec x=\vec b$要么无解（$\vec b$不在像空间中），要么有无穷多解（存在自由变量），不可能有唯一解——而逆矩阵$A^{-1}$存在的前提是方程对任意$\vec b$有唯一解$\vec x=A^{-1}\vec b$，故$A^{-1}$不存在。

6.（开放题，示例）行列式与特征值：对于$A=\begin{pmatrix}2&1\\1&2\end{pmatrix}$，$\det A=3$，特征值$\lambda_1=1,\lambda_2=3$，满足$\lambda_1\lambda_2=1\times3=3=\det A$。几何解释：行列式是面积的总放大倍率，而特征值的乘积给出的是沿各特征方向的"分量放大倍率"之积——两者必然相等，因为面积放大倍率与坐标系选择无关。

\noindent\textbf{课后练习 A组}\\
1. $A^{-1}=\begin{pmatrix}1&-2&1\\0&1&-2\\0&0&1\end{pmatrix}$（逆切变矩阵，每个非对角元取反号）。\\
2. 按第一行展开：$2\cdot\begin{vmatrix}2&1&0\\1&2&1\\0&1&2\end{vmatrix}-1\cdot\begin{vmatrix}1&1&0\\0&2&1\\0&1&2\end{vmatrix}$。第一个三阶$=2(4-1)-1(2-0)=6-2=4$；第二个$=1\cdot(4-1)=3$。总$=2\times4-3=5$。\\
3. $\det A=2\cdot(-3)-1\cdot1=-7$。$A_1=\begin{pmatrix}5&1\\-1&-3\end{pmatrix}$，$\det=-15+1=-14$，$x=2$。$A_2=\begin{pmatrix}2&5\\1&-1\end{pmatrix}$，$\det=-2-5=-7$，$y=1$。\\
4. $3\times4$矩阵不是方阵，没有行列式（行列式仅定义于方阵），也没有特征值（特征值仅定义于方阵）。秩最大为$\min(3,4)=3$。\\
5. $A^{\text T}A$的特征值为$3^2=9$和$(-2)^2=4$。\\
6. 克拉默法则：公式优美，适合小规模（$n\leq3$）和理论推导；计算量巨大（$n!$量级），不适合数值计算。高斯消元：计算效率高（$n^3$量级），适合实际求解；但不如克拉默法则在理论分析中方便。

\noindent\textbf{B组}\\
1. $(A^{\text T}A)^{\text T}=A^{\text T}(A^{\text T})^{\text T}=A^{\text T}A$，故为对称矩阵。对任意向量$\vec x$，$\vec x^{\text T}A^{\text T}A\vec x = (A\vec x)^{\text T}(A\vec x) = \|A\vec x\|^2 \geq 0$，故$A^{\text T}A$是半正定矩阵，特征值均非负。\\
2. 按最后一行拉普拉斯展开：$\det A = a_{nn}\cdot(-1)^{n+n}\det A_{n-1}=a_{nn}\det A_{n-1}$，其中$A_{n-1}$为划去最后一行最后一列的上三角子矩阵（仍为上三角）。递归应用即得$\det A = a_{11}a_{22}\cdots a_{nn}$。\\
3. $A=U\Sigma V^{\text T}=\begin{pmatrix}0&1\\1&0\end{pmatrix}\begin{pmatrix}3&0\\0&1\end{pmatrix}\begin{pmatrix}1&0\\0&1\end{pmatrix}=\begin{pmatrix}0&1\\3&0\end{pmatrix}$。几何上：先不做旋转（$V=I$），沿$x$和$y$分别拉伸$3$倍和$1$倍，然后交换$x$和$y$的坐标。\\
4. 利用行列式性质：将第一行乘以$-x_1$加到第二行，第一行乘以$-x_1$加到第三行；再将新第二行乘以$-(x_2+x_1)$加到第三行，提取公因式即得结果。或用拉普拉斯展开结合归纳法证明。\\
5. $\det A=\begin{vmatrix}2&-3&1\\1&1&-2\\3&-1&1\end{vmatrix}=2(1-2)-(-3)(1+6)+1(-1-3)=-2+21-4=15$。$A_1$：替换第一列，$\det A_1=\begin{vmatrix}1&-3&1\\4&1&-2\\2&-1&1\end{vmatrix}=1(1-2)-(-3)(4+4)+1(-4-2)=-1+24-6=17$，$x=\frac{17}{15}$。$A_2$：$\det A_2=\begin{vmatrix}2&1&1\\1&4&-2\\3&2&1\end{vmatrix}=2(4+4)-1(1+6)+1(2-12)=16-7-10=-1$，$y=-\frac1{15}$。$A_3$：$\det A_3=\begin{vmatrix}2&-3&1\\1&1&4\\3&-1&2\end{vmatrix}=2(2+4)-(-3)(2-12)+1(-1-3)=12-30-4=-22$，$z=-\frac{22}{15}$。\\
6. 特征值为$4,1,-2$，标准形为$4X_1^2+X_2^2-2X_3^2$。不确定（有正有负）。\\
7. 按第一行展开（该行有两个$1$，可优先处理）：$\det=1\cdot(-1)^{1+1}\begin{vmatrix}0&1&0\\1&3&2\\0&2&1\end{vmatrix}+2\cdot(-1)^{1+2}\begin{vmatrix}2&1&0\\0&3&2\\1&2&1\end{vmatrix}+0+1\cdot(-1)^{1+4}\begin{vmatrix}2&0&1\\0&1&3\\1&0&2\end{vmatrix}$。第一个三阶$=0\cdot(3-4)-1\cdot(1-0)+0= -1$。第二个三阶按第一列展开得$2\cdot(3-4)-0+1\cdot(2-0)= -2+2=0$，乘$-2$得$0$。第三个三阶按第一列展开得$2\cdot(2-0)-0+1\cdot(0-1)=4-1=3$，乘$(-1)^{5}=-1$得$-3$。总$\det=1\cdot(-1)+0+(-1)\cdot3=-4$。

\noindent\textbf{7.5 $n$维仿射变换}\\
1. $4\times4$齐次矩阵为$\begin{pmatrix}0&-1&0&5\\1&0&0&0\\0&0&1&0\\0&0&0&1\end{pmatrix}$。\\
2. 齐次坐标将线性变换和平移统一为同一种运算（矩阵乘法），使得任意多次复合都可以简单地通过矩阵连乘完成，无需分别追踪平移分量。这大幅简化了计算机程序的实现。

\noindent\textbf{7.6 $n$维二次型与正定性}\\
1. $M$的对角元全为$2$，非对角元较小（$1$），且矩阵是对称正定的（所有顺序主子式为正），因此三个特征值均应为正。$Q(\vec x)>0$对所有非零$\vec x$成立。\\
2. 二维分类看两个特征值的符号（同号→椭圆，异号→双曲线，有零→抛物线）。$n$维正定性推广了这个思想：看$n$个特征值的符号——全正为正定（$n$维椭球），有正有负为不定（$n$维双曲面），有零为半定（退化）。核心逻辑完全相同：特征值的符号决定了二次型的几何形状。

